\scanpage{25}
\mspage{\handwritten{16}}
définissant des hypersurfaces $H_o,..H_{n-1}$ telles que $X\handwritten{\cap}H_o\handwritten{\cap}..\handwritten{\cap}H_{n-1}=Z$ soit un schéma fini \rlap{\begin{tikzpicture}[baseline=0pt,x=1cm,y=1cm]\draw[line width=.4pt] (0,-.17) .. controls (-.12,.1) and (.18,.28) .. (.70,.36) .. controls (1.08,.43) and (1.18,.52) .. (1.03,.67);\end{tikzpicture}}\inserted{\handwritten{étale}} sur $k$, de degré $rd_1^n$, dont tou\deleted{t}s\deleted{\unsure{es}} les points sont rationnels sur $k$\rlap{\begin{tikzpicture}[baseline=0pt,x=1cm,y=1cm]\draw[line width=.4pt] (0,-.24) .. controls (-.12,.18) and (.03,.38) .. (.39,.39) .. controls (1.0,.38) and (2.75,.37) .. (3.2,.43) .. controls (3.48,.47) and (3.55,.57) .. (3.43,.66);\end{tikzpicture}}\inserted{et distinct des $x_i$.}. \deleted{\unsure{S}}Quitte à augmenter $k$ et $d_1$, on peut supposer qu'il existe une forme $\phi_n$ de même degré $d_1$, définissant $H_n$, tel\handwritten{\unsure{le|l|e}} que $H_n\handwritten{\cap}Z$ soit de cardinal $N$. Pour tout entier $d_2\handwritten{\geq}n+1$, on sait qu'il existe une forme $\underline{\phi}$ à $n+1$ variables \deleted{telle que} définie sur $k$, de degré $d_2$, \deleted{telle que} n'ayant pas de zéro non trivial sur $k$. Alors
\[
\phi=\underline{\phi}\,(\phi_o,\ldots,\phi_n)
\]
est une forme de degré $d_2d_1$ sur $P^r_k$, telle que $(H_{\phi}\handwritten{\cap}X)(k)=(H_n\handwritten{\cap}Z)(k)$ soit de cardinal $N$. De plus, comme cet ensemble ne contient pas les $x_i$, $\phi$ définit bien un point de $U_{d_1d_2}$. Lorsque $N$ est 0 ou 1\rlap{\begin{tikzpicture}[baseline=0pt,x=1cm,y=1cm]\draw[line width=.4pt] (0,-.24) .. controls (-.12,.06) and (-.02,.28) .. (.38,.35);\draw[line width=.25pt,opacity=.55] (.48,.36) .. controls (1.65,.44) and (3.20,.25) .. (4.1,.39);\end{tikzpicture}}\inserted{(plus généralement $N\handwritten{\leq}r$)}, on a terminé, car on peut prendre $d_1=1$, et 7.6 marche avec $d'_o=n+1$. Si on a $N>r$, on note que pour tout $d_{\handwritten{3}}\handwritten{\geq}n+1$, il existe une forme $\unsure{\psi|\mu}$ de degré $d_{\handwritten{3}}$ \rlap{\begin{tikzpicture}[baseline=0pt,x=1cm,y=1cm]\draw[line width=.4pt] (0,-.24) .. controls (-.12,.13) and (.03,.32) .. (.28,.34)--(.62,.32);\draw[line width=.3pt] (.70,.32) .. controls (1.00,.32) and (1.22,.34) .. (1.19,.49);\end{tikzpicture}}\inserted{rat/$k$} telle que $(H_{\psi}\handwritten{\cap}X)(k)=\phi$ d'après ce qu'on vient de voir, et alors la forme $\phi\ \psi$ est de degré \rlap{\begin{tikzpicture}[baseline=0pt,x=1cm,y=1cm]\draw[line width=.4pt] (0,-.2) .. controls (-.04,.08) and (.35,.25) .. (.40,.47) .. controls (.44,.56) and (.39,.64) .. (.36,.70);\end{tikzpicture}}\inserted{$d=$} $d_1d_2+d_{\handwritten{3}}$ et $(H_{\phi\psi}\handwritten{\cap}X)(k)=Z(k)$ a comme cardinal $N$, et $\phi\psi$ définit encore un élément de $U_d$. On peut donc prendre alors $d'_o=d_1\handwritten{\unsure{d|\phi}}_2+(n+1)=(d_1+1)(n+1)$ (en prenant, comme il est loisible, $d_2=n+1$.

\uline{Remarques} 7.7. \handwritten{a)} On trouve ainsi une estimation assez précise da $d'_o$ de 7.6 \deleted{\unsure{(donc}} en fonction de $N$ et de $r$, donc du $d_o$ de 7.1 en fonction de $N=p-1$ et de $r$, \deleted{\unsure{lorsqu'on désire dans}} et du plus petit entier rendant valables les égalités (*). Lorsque dans la conclusion de 7.1 on se contente d'affirmer que les éléments de $G_o$ donnent lieu à au moins deux traces distinctes (ce qui suffit à conclure que $G_o$ n'est pas un $p$-groupe et a fortiori n'est pas le groupe unité, donc $\dim V\handwritten{\geq}1$ i.e. frobenius sur $V$ n'est pas nilpotent) on peut prendre \deleted{pour} simplement $d'_o=n+1$, et pour $d_o$ le plus petit entier $\handwritten{\geq}n+1$ tel que l'on ait (*). Ce choix de $d_o$ est suffisant aussi, par conséquent, pour avoir la validité des
% archive sheet number: 24
